\documentclass[12pt,a4paper]{article}
\usepackage[utf8]{vietnam}
\usepackage{amsmath,amssymb}
\usepackage{geometry}
\geometry{
  top=12mm,
  bottom=12mm,
  left=14mm,
  right=14mm,
  footskip=6mm,
  headheight=14pt
}

\usepackage{tikz}
\usetikzlibrary{arrows.meta,calc,shapes.geometric}
\usepackage[most]{tcolorbox}
\usepackage{fontawesome}
\usepackage{fancyhdr}
\usepackage{enumitem}

\pagestyle{fancy}
\fancyhf{}
\fancyhead[L]{\footnotesize\color{teal!80!black}\faBook\ \textbf{SỔ TAY KIẾN THỨC TOÁN 3}}
\fancyhead[R]{\footnotesize\color{teal!80!black}\textbf{Lớp Toán Thầy Lĩnh -- loptoanthaylinh.vn}}
\fancyfoot[C]{\footnotesize\color{gray}-- Sổ tay kiến thức: Bài 19. Hình tam giác, hình tứ giác, hình chữ nhật, hình vuông --}
\renewcommand{\headrulewidth}{0.4pt}
\renewcommand{\footrulewidth}{0pt}

\setlength{\parindent}{0pt}
\setlength{\parskip}{0pt}

\begin{document}

% --- BANNER BÀI HỌC ---
\begin{tcolorbox}[
  enhanced,
  colback=teal!8!white,
  colframe=teal!70!black,
  arc=3mm, boxrule=1.2pt,
  left=3mm, right=3mm, top=2mm, bottom=2mm,
  halign=center
]
  {\LARGE\bfseries\color{teal!90!black}\faBookmark\ BÀI 19: HÌNH TAM GIÁC, TỨ GIÁC, CHỮ NHẬT, VUÔNG}\\[1mm]
  {\itshape\footnotesize\color{black!75} Nhận dạng số đỉnh, số cạnh, đặc điểm góc vuông và cạnh bằng nhau của các hình phẳng trên 1 trang A4}
\end{tcolorbox}

\vspace{-2mm}

\noindent
\begin{minipage}[t]{0.49\textwidth}\vspace{0pt}
\begin{tcolorbox}[
  enhanced,
  colback=blue!4!white,
  colframe=blue!75!black,
  arc=2mm, boxrule=0.8pt,
  left=2mm, right=2mm, top=1.8mm, bottom=1.8mm,
  title={\footnotesize\bfseries\color{white}\faSquareO\ 1. Hình tam giác \& Hình tứ giác}
]
{\footnotesize
\begin{itemize}\setlength{\itemsep}{1pt}\setlength{\leftskip}{-2.5mm}
  \item \textbf{Hình tam giác:}
  \begin{itemize}\setlength{\itemsep}{0.5pt}\setlength{\leftskip}{-2.5mm}
    \item Có $3$ đỉnh, $3$ cạnh và $3$ góc.
    \item Ví dụ tam giác $ABC$: Đỉnh $A, B, C$; Cạnh $AB, BC, CA$.
  \end{itemize}
  \item \textbf{Hình tứ giác:}
  \begin{itemize}\setlength{\itemsep}{0.5pt}\setlength{\leftskip}{-2.5mm}
    \item Có $4$ đỉnh, $4$ cạnh và $4$ góc.
    \item Ví dụ tứ giác $MNPQ$: Đỉnh $M, N, P, Q$; Cạnh $MN, NP, PQ, QM$.
  \end{itemize}
\end{itemize}
}
\end{tcolorbox}
\end{minipage}
\hfill
\begin{minipage}[t]{0.49\textwidth}\vspace{0pt}
\begin{tcolorbox}[
  enhanced,
  colback=teal!4!white,
  colframe=teal!75!black,
  arc=2mm, boxrule=0.8pt,
  left=2mm, right=2mm, top=1.8mm, bottom=1.8mm,
  title={\footnotesize\bfseries\color{white}\faSquareO\ 2. Hình chữ nhật \& Hình vuông}
]
{\footnotesize
\begin{itemize}\setlength{\itemsep}{1pt}\setlength{\leftskip}{-2.5mm}
  \item \textbf{Hình chữ nhật (Tứ giác đặc biệt):}
  \begin{itemize}\setlength{\itemsep}{0.5pt}\setlength{\leftskip}{-2.5mm}
    \item Có $4$ góc vuông.
    \item $2$ cạnh dài bằng nhau ($a$) và $2$ cạnh ngắn bằng nhau ($b$).
  \end{itemize}
  \item \textbf{Hình vuông (Hình chữ nhật hoàn hảo):}
  \begin{itemize}\setlength{\itemsep}{0.5pt}\setlength{\leftskip}{-2.5mm}
    \item Có $4$ góc vuông.
    \item Có $\mathbf{4}$ \textbf{cạnh dài bằng nhau}.
  \end{itemize}
  \item Hình vuông là một dạng hình chữ nhật đặc biệt.
\end{itemize}
}
\end{tcolorbox}
\end{minipage}

\vspace{-2.5mm}
% --- KHỐI 3: SƠ ĐỒ TƯ DUY TIKZ ---
\noindent
\begin{tcolorbox}[
  enhanced,
  colback=orange!4!white,
  colframe=orange!80!black,
  arc=2mm, boxrule=0.8pt,
  left=2mm, right=2mm, top=1.5mm, bottom=1.5mm,
  title={\footnotesize\bfseries\color{white}\faSitemap\ SƠ ĐỒ TƯ DUY: HỆ QUAN HỆ 4 HÌNH PHẲNG CỐT LÕI}
]
\centering
\begin{tikzpicture}[
  fig/.style={draw=blue!80!black, fill=blue!10, rounded corners=1pt, minimum height=0.7cm, font=\scriptsize\bfseries}
]
  \node[fig, text width=2.4cm, align=center] (t1) at (0, 0) {TAM GIÁC\\$3$ đỉnh, $3$ cạnh\\$3$ góc};
  \node[fig, text width=2.4cm, align=center, fill=purple!12, draw=purple!80!black] (t2) at (3.5, 0) {TỨ GIÁC\\$4$ đỉnh, $4$ cạnh\\$4$ góc};
  \node[fig, text width=2.5cm, align=center, fill=teal!15, draw=teal!80!black] (t3) at (7.0, 0) {HÌNH CHỮ NHẬT\\$4$ góc vuông\\Cặp cạnh đối bằng nhau};
  \node[fig, text width=2.4cm, align=center, fill=orange!15, draw=orange!80!black] (t4) at (10.4, 0) {HÌNH VUÔNG\\$4$ góc vuông\\$\mathbf{4}$ cạnh bằng nhau};

  \draw[thick, ->, purple!80!black] (t2) -- (t3) node[midway, above, font=\tiny\bfseries] {Có $4$ góc vuông};
  \draw[thick, ->, orange!80!black] (t3) -- (t4) node[midway, above, font=\tiny\bfseries] {Dài = Rộng};
\end{tikzpicture}
\end{tcolorbox}

\vspace{-2.5mm}
% --- KHỐI 4: VÍ DỤ & MẸO TRÁNH BẪY ---
\noindent
\begin{minipage}[t]{0.49\textwidth}\vspace{0pt}
\begin{tcolorbox}[
  enhanced,
  colback=green!4!white,
  colframe=green!60!black,
  arc=2mm, boxrule=0.8pt,
  left=2mm, right=2mm, top=1.8mm, bottom=1.8mm,
  title={\footnotesize\bfseries\color{white}\faCheckCircle\ Ví dụ thực hành chuẩn mực}
]
{\footnotesize
\textbf{1. Phân loại hình:}\\
Tứ giác có $4$ góc vuông và $4$ cạnh đều dài $5\text{ cm}$ là: \textbf{Hình vuông}.\\
Tứ giác có $4$ góc vuông, chiều dài $8\text{ cm}$, chiều rộng $5\text{ cm}$ là: \textbf{Hình chữ nhật}.\\
\textbf{2. Đếm hình ghép:}\\
Hình vuông bị chia bởi $1$ đường chéo sẽ tạo thành: $\mathbf{2}$ \textbf{hình tam giác}.
}
\end{tcolorbox}
\end{minipage}
\hfill
\begin{minipage}[t]{0.49\textwidth}\vspace{0pt}
\begin{tcolorbox}[
  enhanced,
  colback=red!4!white,
  colframe=red!70!black,
  arc=2mm, boxrule=0.8pt,
  left=2mm, right=2mm, top=1.8mm, bottom=1.8mm,
  title={\footnotesize\bfseries\color{white}\faExclamationTriangle\ Mẹo tránh bẫy trắc nghiệm}
]
{\footnotesize
\begin{itemize}\setlength{\itemsep}{1.5pt}\setlength{\leftskip}{-2.5mm}
  \item \textbf{Bẫy chỉ nhìn cạnh mà quên góc:}\\
  Tứ giác có 4 cạnh bằng nhau chưa chắc là hình vuông nếu 4 góc của nó không phải là góc vuông (như hình thoi)!
  \item \textbf{Khẳng định chuẩn:}\\
  "Mọi hình vuông đều là hình chữ nhật" là ĐÚNG. Nhưng "Mọi hình chữ nhật đều là hình vuông" là SAI!
\end{itemize}
}
\end{tcolorbox}
\end{minipage}


\end{document}
